\section{Formalisierung in Lean 4}

Die Datei \texttt{EllipseSmoothing.lean} bildet den bisherigen Beweisstand ab.
Der geometrische Kern wird divisionsfrei formuliert:

\begin{lstlisting}
def OnEllipse (E : Ellipse) (p : R × R) : Prop :=
  E.b ^ 2 * (p.1 - E.cx) ^ 2 +
    E.a ^ 2 * (p.2 - E.cy) ^ 2 =
      E.a ^ 2 * E.b ^ 2
\end{lstlisting}

Der zentrale Satz abstrahiert vollständig von der Herkunft des Kontrollwertes:

\begin{lstlisting}
theorem ellipseLift_onEllipse
    (E : Ellipse) {g branch : R}
    (hg : g ^ 2 ≤ 1)
    (hbranch : branch ^ 2 = 1) :
    OnEllipse E (ellipseLift E g branch)
\end{lstlisting}

Bereits formal bewiesen sind:
\begin{itemize}
  \item die Ellipseninvariante des allgemeinen Wurzellifts,
  \item die Unabhängigkeit vom oberen oder unteren Zweig,
  \item die Endpunkte und Beschränktheit der linearen Dreieckssteuerung,
  \item die Vererbung der Invariante an jede zulässige Steuerfunktion,
  \item die Summierbarkeit und Nichtnegativität der Glättungsgewichte für
        $0\leq\lambda\leq1$,
  \item die scharfe untere Schranke $1\leq D_\lambda$ für den Nenner,
  \item die globale Schranke $\lvert G_\lambda(\theta)\rvert\leq1$,
  \item die Ellipseninvariante der geglätteten Bahn im gesamten
        abgeschlossenen Glättungsintervall.
\end{itemize}

Der Beweis benutzt die konvergente Majorante $(2k+1)^{-2}$. Weil alle Gewichte
im Intervall $0\leq\lambda\leq1$ nichtnegativ sind und der Summand für $k=0$
genau eins ist, gilt $D_\lambda\geq1$. Aus
$\lvert\cos((2k+1)\theta)\rvert\leq1$ folgt dann
$\lvert N_\lambda(\theta)\rvert\leq D_\lambda$ und nach Division die globale
Zulässigkeit.

Noch offen sind die exakte Identifikation der beiden Randfälle und die
Glattheit für $0<\lambda$. Auch diese Aussagen werden nicht durch zusätzliche
Axiome ersetzt, sondern bleiben explizite Ziele des nächsten Gates.
